\section{The ordered Poincar\'e--Birkhoff--Witt theorem, enveloping Hopf algebras, and the completed tensor product}\label{app:pbw}
This appendix proves the injectivity statement used in \cref{sec:universal},
constructs the Hopf structure on an arbitrary enveloping algebra, and shows
concretely why the coproduct of an infinite series needs a completed tensor
product.

\subsection{The universal enveloping algebra and its ordered normal forms}
Let $\h$ be a Lie algebra over $\kk$. Its universal enveloping algebra is
\begin{equation}\label{eq:Udef}
 U(\h)=T(\h)\big/\bigl\langle u\otimes v-v\otimes u-[u,v]:u,v\in\h\bigr\rangle,
\end{equation}
the quotient of the tensor algebra on the vector space $\h$ by the
two-sided ideal generated by the displayed relations. A linear map from
$\h$ to an associative algebra $A$ extends uniquely to an algebra
homomorphism on $T(\h)$, and this extension annihilates the ideal exactly
when the map preserves brackets (with $A$ given its commutator bracket);
this is the universal property. It remains to see that $\h\to U(\h)$ is
injective.

\begin{theorem}[Ordered PBW theorem]\label{thm:pbw}
Let $(e_i)_{i\in I}$ be a totally ordered basis of $\h$. The ordered
monomials $e_{i_1}e_{i_2}\cdots e_{i_m}$ with $i_1\leq i_2\leq\cdots\leq i_m$,
together with the empty monomial $1$, form a vector-space basis of $U(\h)$.
In particular the natural map $\h\to U(\h)$ is injective.
\end{theorem}
\begin{proof}
Work in $T(\h)$ with the basis of words in the letters $e_i$. Consider the
rewriting rule, applicable to any adjacent inverted pair $j>i$ inside a
word,
\begin{equation}\label{eq:rewrite}
 \cdots e_je_i\cdots\ \longmapsto\ \cdots e_ie_j\cdots+\cdots[e_j,e_i]\cdots,
\end{equation}
where $[e_j,e_i]$ is expanded as a finite linear combination of basis
letters. Assign to a word the pair (length, number of inversions), ordered
lexicographically with length first. The swapped word has the same length
and one inversion fewer, and every word in the bracket term is shorter. So
each application of the rule replaces a word by a linear combination of
words of strictly smaller measure. Since the measure takes values in a
well-ordered set, every sequence of rewritings of a fixed word terminates,
and the result is a linear combination of ordered words (words with no
inverted adjacent pair). This proves that ordered words span $U(\h)$.

We prove that the final result does not depend on the choices made
(confluence), by induction on the measure of the starting word $w$. If two
different first steps are available for $w$ at disjoint adjacent pairs,
performing the other one next in either case leads to the same linear
combination, whose words all have smaller measure than $w$; by the
induction hypothesis each of those words has a unique normal form, hence
so does the combination, and both first choices lead to it. If the two
available first steps overlap, the overlap is a factor $e_ke_je_i$ with
$k>j>i$ inside $w=p\,e_ke_je_i\,q$. Reducing the left pair first and then
continuing to order the three original letters gives, before any
normalization of the shorter terms,
\[
 p\bigl(e_ie_je_k+[e_j,e_i]e_k+e_j[e_k,e_i]+[e_k,e_j]e_i\bigr)q,
\]
and reducing the right pair first gives
\[
 p\bigl(e_ie_je_k+e_i[e_k,e_j]+[e_k,e_i]e_j+e_k[e_j,e_i]\bigr)q.
\]
(Each line records the sequence of swaps that brings $e_ke_je_i$ to
$e_ie_je_k$, collecting the bracket terms produced.) Their difference is
$p\,\Xi\,q$ with
\[
 \Xi=[e_j,e_i]e_k-e_k[e_j,e_i]+e_j[e_k,e_i]-[e_k,e_i]e_j+[e_k,e_j]e_i-e_i[e_k,e_j].
\]
Every word in $\Xi$ has length two, hence smaller measure than $e_ke_je_i$,
and the same holds after attaching $p$ and $q$. Applying the rule to each
length-two product $ab-ba$ with $a,b\in\h$ (which, expanded in the basis and
reduced, has normal form equal to the normal form of $[a,b]$; this is the
content of \eqref{eq:rewrite} for a single pair, applied bilinearly)
shows that $p\,\Xi\,q$ has the same normal form as
\[
 p\bigl([[e_j,e_i],e_k]+[e_j,[e_k,e_i]]+[[e_k,e_j],e_i]\bigr)q
 =p\bigl([[e_j,e_i],e_k]+[[e_i,e_k],e_j]+[[e_k,e_j],e_i]\bigr)q=0,
\]
by antisymmetry and the Jacobi identity in $\h$, where all normal forms of
the shorter words involved are unique by the induction hypothesis. Hence
the two first choices lead to combinations with the same normal form.
Since every rule has length two, there are no other overlapping patterns.
Thus any two complete reductions of $w$ agree: compare their first steps
as above, then apply the induction hypothesis to all words below $w$.
Extending by linearity yields a well-defined \emph{normal-form map}
$\nu:T(\h)\to T(\h)$ whose image is the span of ordered words and which
fixes ordered words.

Each rewriting step changes an element by an element of the ideal
$J$ in \eqref{eq:Udef}: it replaces $p\,e_je_i\,q$ by
$p(e_ie_j+[e_j,e_i])q$, and the difference is $p(e_je_i-e_ie_j-[e_j,e_i])q\in J$.
Hence $x-\nu(x)\in J$ for every $x$. Conversely, every generator
$r=uv-vu-[u,v]$ of $J$, multiplied on either side by words, has normal form
zero: by bilinearity it suffices to take $u=e_j$, $v=e_i$, and for $j>i$
the two sides $e_je_i$ and $e_ie_j+[e_j,e_i]$ of the rule have the same
normal form by construction, while for $j=i$ the relation is $0$ and for
$j<i$ it is the negative of a rule instance. Thus $\nu(J)=0$, and
$\nu$ descends to a linear map $U(\h)=T(\h)/J\to\mathrm{span}(\text{ordered words})$
that is surjective (ordered words are fixed) and injective (if $\nu(x)=0$
then $x=x-\nu(x)\in J$). Therefore the ordered words form a basis of
$U(\h)$. Words of length one are ordered, so the basis vectors $e_i$ remain
linearly independent in $U(\h)$, which proves injectivity of $\h\to U(\h)$.
\end{proof}
The proof uses only bilinearity, antisymmetry, and the Jacobi identity; it
is independent of the BCH theorem and of any matrix realization, so its use
in \cref{sec:universal} introduces no circularity.

\subsection{The free Lie algebra and the free associative algebra}
Let $F(X,Y)$ be the abstract free Lie algebra on two generators, the
quotient of the free nonassociative algebra on bracketed words by the
ideal generated by antisymmetry and Jacobi. A Lie homomorphism out of
$F(X,Y)$ is the same as a choice of two elements of the target: substitution
defines the map, and the imposed relations are precisely what makes it well
defined. Hence the universal property of $U(F(X,Y))$ says that an
associative homomorphism from it to any associative algebra is exactly a
choice of two elements, which is also the universal property of the free
associative algebra $\kk\langle X,Y\rangle=T(V)$; the mutually inverse
homomorphisms exchanging the generators identify the two, as in the proof
of \cref{prop:enveloping}. By \cref{thm:pbw}, $F(X,Y)$ injects into
$U(F(X,Y))\cong T(V)$, and its image is the Lie subalgebra $\cL(X,Y)$
generated by $X,Y$. This identifies the concrete $\cL(X,Y)$ used throughout
the article with the abstract free Lie algebra, and shows that an equality
of bracket polynomials verified as an equality of words in $T(V)$ is an
identity in every Lie algebra.

\subsection{The enveloping Hopf structure}
On the tensor algebra $T(\h)$ set, for $u\in\h$,
\begin{equation}\label{eq:UHopf}
 \Delta u=u\otimes1+1\otimes u,\qquad\eps(u)=0,\qquad S(u)=-u,
\end{equation}
and extend $\Delta,\eps$ as algebra homomorphisms and $S$ as an
antihomomorphism. For a defining relation $r=uv-vu-[u,v]$ of
\eqref{eq:Udef}, direct expansion gives
\[
 \Delta r=(uv-vu)\otimes1+1\otimes(uv-vu)-[u,v]\otimes1-1\otimes[u,v]
 =r\otimes1+1\otimes r,
\]
because the mixed terms $u\otimes v+v\otimes u-v\otimes u-u\otimes v$ cancel;
also $\eps(r)=0$ and $S(r)=vu-uv+[u,v]=-r$. Hence $\Delta(J)\subseteq
J\otimes T(\h)+T(\h)\otimes J$, $\eps(J)=0$, and $S(J)\subseteq J$, so the
three maps descend to $U(\h)$. Coassociativity and the counit identities
hold on generators, hence on all of $U(\h)$ by multiplicativity; the
antipode identities $m(S\otimes\id)\Delta=\eta\eps=m(\id\otimes S)\Delta$
hold on generators ($-u+u=0=\eps(u)$) and propagate to products by the
inductive argument of \cref{lem:antipode}, which used only that $\Delta$
is multiplicative and $S$ antimultiplicative. Thus $U(\h)$ is a Hopf
algebra, and under $U(\cL(X,Y))\cong T(V)$ this is exactly the unshuffle
Hopf algebra of \cref{sec:hopf}; degree completion gives its formal
power-series version. For a general $\h$, tensor length is not a grading
of $U(\h)$, since a bracket relation changes length; formal exponentials
are then taken in $U(\h)[[t]]$ with arguments $tx,ty$, which preserves every
coefficient identity while avoiding an inappropriate grading.

\subsection{Why the coproduct needs a completed tensor product}
Restrict to the one-generator subalgebra and consider
$F=\sum_{n\geq0}X^n$. By \eqref{eq:unshuffle},
\[
 \coeff{X^p\otimes X^q}\Delta F=\binom{p+q}{p}\qquad(p,q\geq0),
\]
since $X^{p+q}$ has $\binom{p+q}p$ subsets of positions of size $p$. The
leading $(N+1)\times(N+1)$ block of this coefficient matrix is
$C=(\binom{p+q}p)_{0\leq p,q\leq N}$. By Vandermonde's identity,
$\binom{p+q}p=\sum_k\binom pk\binom qk$ (compare coefficients of $t^p$ in
$(1+t)^p(1+t)^q=(1+t)^{p+q}$, using $\binom{q}{p-k}=\binom q{q-p+k}$ and
reindexing), so $C=PP^{\mathsf T}$ with $P_{pk}=\binom pk$ for $k\leq p$ and
$P_{pk}=0$ otherwise. Since $P$ is lower triangular with diagonal entries
$1$, $\det C=(\det P)^2=1$, so $C$ has rank $N+1$ for every $N$. On the
other hand, if $\Delta F$ were a finite sum $\sum_{j=1}^mA_j\otimes B_j$ of
separated tensors with $A_j,B_j\in\kk[[X]]$, its coefficient matrix would
be $\sum_j\alpha_j\beta_j^{\mathsf T}$ with $\alpha_j,\beta_j$ the
coefficient vectors of $A_j,B_j$, of rank at most $m$. Taking $N\geq m$
gives a contradiction. Thus $\Delta F$ does not lie in the algebraic tensor
product $\widehat T\otimes\widehat T$, while it does lie in the completed
tensor product used in \cref{sec:hopf}, whose elements are arbitrary
families of coefficients indexed by pairs of words.
